\chapter{KL 估计量、无偏性与方差}

\section{K4 的无偏性}

固定 history，记 $p=\pi_\theta(\cdot\mid s)$、$q=\mu_k(\cdot\mid s)$，并从 $p$ 采样 $a$。定义
\[
 w(a)=\frac{q(a)}{p(a)},
 \qquad r_\theta(a)=\frac{p(a)}{\sg(p(a))}.
\]
数值上 $r_\theta=1$，但其导数为 $\nabla\log p(a)$。

\begin{theorem}[K4 value and gradient]
\label{thm:k4}
若 $p\ll q$ 且二阶矩有限，则
\[
 \E_{a\sim p}[r_\theta(a)(-\log w(a))]=\KL(p\Vert q),
\]
并且自动微分意义下
\[
 \E_{a\sim p}[\nabla_\theta\{r_\theta(a)(-\log w(a))\}]
 =\nabla_\theta\KL(p\Vert q).
\]
\end{theorem}

\begin{proof}
数值 value 中 $r_\theta=1$，所以第一式是
\[
\E_p[-\log(q/p)]=\sum_a p(a)\log(p(a)/q(a)).
\]
对 stop-gradient 构造求导时，$-\log w=\log p-\log q$ 的 derivative 与 $r_\theta$ 的 derivative 共同给出
\[
 \nabla_\theta\{r_\theta(-\log w)\}
 =\nabla_\theta\log p(a)\,[\log p(a)-\log q(a)]
 +\nabla_\theta\log p(a).
\]
最后一项在 $p$ 下期望为零，第一项正是 KL 的 score-function derivative。\hfill$\square$
\end{proof}

K1 的数值也无偏，但其 autograd gradient 为零；K3 的 gradient 对应 $\KL(q\Vert p)$ 的方向；K2 梯度正确但 value 有偏。上述构造针对 MDPO 中的 $\KL(p\Vert q)$。BPO 的 student center 是相反方向的 $\KL(q\Vert p)$，应在冻结 anchor $q$ 下估计 $\log(q/p)$；PMD-Mean 不含这个 student KL center。

\section{固定 anchor 测度上的重要性校正}

若工程上必须从 $q=\mu_k$ 采样，令 $\rho_\theta(a)=p(a)/q(a)$，并在积分权重上停止梯度。对于任意可积向量 $h$，
\[
 \E_q[\sg(\rho_\theta)h(a)]=\E_p[h(a)].
\]
因此 MDPO 的固定 anchor 测度梯度可直接写为
\[
\widehat g_i^{\mathrm{MDPO}}
=\sg\!\left(\frac{p(a_i)}{q(a_i)}\right)
\sg\!\left[f(a_i)-\frac1\eta\log\frac{p(a_i)}{q(a_i)}\right]
\nabla_\theta\log p(a_i),\qquad a_i\sim q.
\]
reward、log-ratio 与 score 的乘积二阶可积时，它无偏；若该乘积范数由 $C$ 上界，则
\begin{equation}\label{eq:bpo-variance}
\mathbb E_q\|\widehat g_i^{\mathrm{MDPO}}\|^2
\le C^2\mathbb E_q[(p/q)^2]
=C^2\bigl(1+\chi^2(p\Vert q)\bigr).
\end{equation}
这里的 $p/q$ 测度校正属于 MDPO；不能据此为 PMD-Mean 或 BPO 添加 student importance ratio。它们原本就在冻结 anchor 下定义，action 从 $q$ 采样时无需这种校正。

对 BPO，完整 conditional center 已知时的实际下降方向为
\[
-\nabla\mathcal L_{\mathrm{BPO}}
=\mathbb E_q\left[
\left(f-\mathbb E_q[f]-\frac{U-\mathbb E_q[U]}\eta\right)
\nabla\log p\right],\qquad U=\log(p/q).
\]
对 PMD-Mean，则为
\[
-\nabla\mathcal L_{\mathrm{Mean}}
=\mathbb E_q\left[
\bigl(\eta(f-\mathbb E_q[f])-U\bigr)\nabla\log p\right].
\]
估计 center 时须区分解析 policy-gradient 权重与 sampled 原始平方损失：前者使用停止梯度的权重，后者的 BPO student center 保留梯度，并采用第二章给出的有限样本系数校正。单个 center 估计量的无偏性不自动保证它嵌入平方损失后的梯度无偏。

\section{Top-$K$ 与完整词表的关系}

如果保存当前策略和 anchor 的 top-$K$ logits，完整 reverse KL 可以拆成
\[
 \sum_{a\in Q_K}p(a)\log\frac{p(a)}{q(a)}
 +\sum_{a\notin Q_K}p(a)\log\frac{p(a)}{q(a)}.
\]
第一项可精确计算；第二项不能被简单删除，否则得到截断 KL 而非完整 KL。2602.04417 的 Top-$K$ 估计量在尾部使用 masked sampled K4：若采样 token 不在 $Q_K$，加入
\[
 \mathbf 1(a\notin Q_K)\,\sg\!\left(\frac{p(a)}{\pi_{\rm sampling}(a)}\right)
 r_\theta(a)(-\log w(a)).
\]
它同时保持 value 和 gradient 无偏。对本项目而言，第一版小实验不使用 top-$K$ reward，而先验证 full-vocabulary oracle 和 sampled K4 两条路径；否则无法判断误差来自 mirror target 还是 tail estimator。

\section{有限 inner 更新的误差}

设方法自己的总体最优分布为 $p_k^{*,M}$，inner 算法在固定 $s$ 上得到 $p_{k,m}$。PMD-Mean 使用 Lambert-$W$ target，另外三种方法使用 Gibbs target $p_k^+$。定义优化误差
\[
 e_{k,m}(s)=\frac12\|p_{k,m}(\cdot\mid s)-p_k^{*,M}(\cdot\mid s)\|_1.
\]
若 loss gradient 是 $L$-smooth，且使用 unbiased gradient with variance at most $\sigma^2/B$，标准 projected SGD 给出
\[
\frac1M\sum_{m=0}^{M-1}\E\|\nabla\mathcal L(\theta_m)\|^2
\leq
\frac{2(\mathcal L_0-\mathcal L_*)}{\alpha M}
 +\frac{L\alpha\sigma^2}{B}.
\]
对采用 student/anchor importance ratio 的 MDPO estimator，方差界包含 $1+\chi^2(p_{k,m}\Vert\mu_k)$。PMD-Mean 和 BPO 的 anchor-measure estimator 应直接分析其 residual、score 与 center 的联合二阶矩，不能套用这个 importance-ratio 结论。

PMD-Mean 到 Gibbs 分布的差异还包含总体目标偏差：
\[
\TV(p_{k,m},p_k^+)
\le \TV(p_{k,m},p_k^{*,\mathrm{Mean}})
+\TV(p_k^{*,\mathrm{Mean}},p_k^+).
\]
增加 inner steps 只能降低第一项，第二项由均值近似决定。
